\section*{第 \ref{ch:g1:my-body} 章——我的身体}

\begin{solution}{exo:g1:my-body:1}
头、躯干、\omterm{def:g1:my-body:parts}{四肢}（两条手臂和两条腿）。颈把
头与躯干连在一起。
\end{solution}

\begin{solution}{exo:g1:my-body:2}
手臂的末端是手；腿的末端是脚。
\end{solution}

\begin{solution}{exo:g1:my-body:3}
(a) 肘。(b) 膝。(c) 腕。
\end{solution}

\begin{solution}{exo:g1:my-body:4}
肩把手臂横向抬起，肘使手臂弯曲，
腕把手歪向耳朵。
\end{solution}

\begin{solution}{exo:g1:my-body:5}
相同之处：整体结构相同——都有头、身体（躯干）、
\omterm{def:g1:my-body:parts}{四肢}。不同之处：你用两条腿直立，两只手臂可以自由活动，
而猫用四条腿行走；猫有尾巴，
你没有。
\end{solution}

\begin{solution}{exo:g1:my-body:6}
只凭眼睛看，你要斜着估计头顶的位置，记号
可能偏高或偏低。把平放的书水平贴紧墙面，
让书的上缘正对头顶并直接贴到墙上：
每次测量都公平、一致。
\end{solution}

\begin{solution}{exo:g1:my-body:7}
它们表明，两个日期之间你的身体长高了。不——
生长太慢，感觉不到；只有记号才能让它显现出来。
\end{solution}

\begin{solution}{exo:g1:my-body:8}
开放性答案：走路、坐下、捡东西、
拍手、写字——几乎每个动作都要用到
膝、髋、肘、肩或腕。没有\omterm{def:g1:my-body:joint}{关节}，身体
就会像一根僵硬的木棍那样运动。
\end{solution}

\begin{solution}{exo:g1:my-body:9}
每个人在自己的身高和速度上都属于正常：有的
孩子长得早，有的长得晚，最终都会达到各自的
成年身高。身高记号是用来观察自己的生长的，
不是用来和别人比赛的。
\end{solution}

\begin{solution}{exo:g1:my-body:10}
在髋处切开（这样腿才能向前弯曲坐下），在
膝处切开（这样小腿可以自然下垂）；
再在颈和肩处切开，便于前倾和保持平衡。切口
模仿了身体的\omterm{def:g1:my-body:joint}{关节}：髋、膝、颈、肩。
\end{solution}
